\documentclass[a4paper,12pt]{article}
\usepackage[norwegian]{babel}
\usepackage{amsmath, amssymb}
\usepackage{geometry}
\usepackage{fancyhdr}
\usepackage{listings}
\usepackage{xcolor}
\usepackage{enumitem}
\newcommand{\D}{\mathrm{d}}
\usepackage[colorlinks=true, linkcolor=red!50!black, citecolor=red!50!black, urlcolor=blue]{hyperref}

\geometry{margin=2.5cm}
\pagestyle{fancy}
\fancyhf{}
\rhead{R2}
\lhead{Parameterframstilling og vektorfunksjoner}
\rfoot{\thepage}

\usepackage{setspace}
\onehalfspacing

\lstset{
    basicstyle=\ttfamily\small,
    backgroundcolor=\color{gray!10},
    frame=single,
    breaklines=true,
    columns=fullflexible,
    keepspaces=true,
    showstringspaces=false,
    upquote=true
}

\title{Parameterframstilling og vektorfunksjoner}
\author{}
\date{}

\begin{document}

\maketitle

\subsection*{Oppgave 1: Parameterframstilling}

Det er flere måter å beskrive bevegelse i flere dimensjoner. Den første vi skal se på er \textit{parameterframstilling}. La

\[
l:
\begin{cases}
x(t)=4t - 2 \\
y(t)=4t+2 \\
z(t)= -2t +6
\end{cases},\, 0 \leq t \leq 5.
\]
Vi sier at linja $l$ er gitt ved $x(t)$, $y(t)$ og $z(t)$ for $t$ mellom 0 og 5. \textit{Parameteren} $t$ er ofte tid. Når $t$ ikke er tid kan du late som.

\begin{enumerate}[label=\alph*)]
\item For å bruke parameterframstilling i GeoGebra må du bruke denne funksjonen:
\begin{lstlisting}
Kurve(<Uttrykk>,<Uttrykk>,<Uttrykk>,<Parametervariabel>,<Startverdi>,<Sluttverdi>)
\end{lstlisting}
Skriv følgende i GeoGebra:
\begin{lstlisting}
l = Kurve(4t-2, 4t+2, -2t+6, t, 0, 5)
\end{lstlisting}

\item Hva representerer endepunktene på grafen til $l$?
\end{enumerate}


\subsection*{Oppgave 2: Punkt}

Vi kan også beskrive bevegelsen som et punkt $A$ med koordinater $(x(t),y(t),z(t))$. I dette tilfellet blir koordinatene $(4t-2,4t+2,-2t+6)$.

\begin{enumerate}[label=\alph*)]
\item Legg til en \lstinline{Glider} for $t$. Skriv følgende i GeoGebra:
\begin{lstlisting}
t = 0
\end{lstlisting}
Høyereklikk på \lstinline{t} og velg \lstinline{Lag glider}.

Høyereklikk på \lstinline{t} igjen og velg \lstinline{Innstillinger}. Gå til \lstinline{Glider} og sett \lstinline{Min} til 0, \lstinline{Maks} til 5, og \lstinline{Animasjonstrinn} til 0.01.

\item Legg til punktet $A$, skriv:
\begin{lstlisting}
A = l(t)
\end{lstlisting}
Dette har samme effekt som å skrive \lstinline{A = (4t-2,4t+2,-2t+6)}. Fordelen er at punktet \lstinline{A} nå oppdateres når linja \lstinline{l(t)} oppdateres.
\item Animer $t$ ved å trykke på startknappen. Hva slags bevegelse observerer du?
\item Hvor langt har punktet $A$ beveget seg fra $t=0$ til $t=5$?\footnote{$\Delta s = 30$.}
\item Hva er hastigheten?\footnote{$v = 6$.}
\end{enumerate}

\subsection*{Oppgave 3: Vektorfunksjon}

Den siste\footnote{Men kanskje også beste.} måten å beskrive bevegelsen på er med en \textit{vektorfunksjon}. Denne får navnet posisjonsvektor (ofte $\vec r(t)$), og er gitt ved
\[\vec r(t) = \overrightarrow{OA}.\]
Med andre ord er posisjonsvektoren $\vec r (t)$ en vektor som peker på punktet der det er akkurat nå. Vi bruker vanlig vektornotasjon:
\[
\vec r(t) = 
\begin{pmatrix}
x(t) \\
y(t) \\
z(t)
\end{pmatrix} =
\begin{pmatrix}
4t-2 \\
4t+2 \\
-2t+6
\end{pmatrix}.
\]

\begin{enumerate}[label=\alph*)]
\item Legg til $\vec r (t)$ i GeoGebra. Skriv følgende:
\begin{lstlisting}
r(t) = Vektor(A)
\end{lstlisting}

\item Fartsvektoren $\vec v(t)$ er gitt som den deriverte av posisjonsvektoren $\vec r(t)$:
\[ \vec v(t) = \frac{d\vec r}{dt}. 
\]
Vis / argumenter for / velg å tro at \[ \vec v(t) = \begin{pmatrix}
4 \\
4 \\
-2
\end{pmatrix}.\]

\item Det går selvsagt også an å gjøre dette i GeoGebra. Skriv følgende:
\begin{lstlisting}
b = Derivert(l)
v = Vektor(A, A + b(t))
\end{lstlisting}
Det kan være en idé å endre farge på $\vec v$. Det er også lurt å skjule $b(t)$ siden denne kun er en hjelpefunksjon for å nå $\vec v$.\newline

Finn hastigheten $|\vec v(t)|$ og sammenlign med 2 e).

\item Bruk samme strategi for akselerasjonsvektoren $\vec a(t)$. Altså skriv:
\begin{lstlisting}
c = Derivert(b)
a = Vektor(A, A + c(t))
\end{lstlisting}

Hvorfor synes ikke $\vec a$?

\end{enumerate}


\section*{Generell bevegelse}

Siden du har vært flink pike/dreng og gjort som fortalt, kan du nå enkelt studere andre bevegelser. La f.eks.

\[
l:
\begin{cases}
x(t)=\cos(2t) \\
y(t)=\sin(2t) \\
z(t) = t - 3
\end{cases},\, 0\leq t \leq 5.
\]


\begin{enumerate}[label=\alph*)]
\item Alt du må gjøre er å endre på definisjonen av $l$ slik at det nå står:
\begin{lstlisting}
l = Kurve(cos(2t), sin(2 t), t-3, t, 0, 5)
\end{lstlisting}

\item Hva blir utrykket for posisjonsvektoren $\vec r(t)$?
\item I hvilket punkt krysses $xy$-planet?\footnote{$(\cos(6),\sin(6),0) \approx (0.96,-0.28,0)$.}
\item Vis at fartsvektoren $\vec v(t)$ er gitt ved
\[
\vec v(t) =
\begin{pmatrix}
-2\sin(2t) \\
2\cos(2t) \\
1
\end{pmatrix}.
\]
\item Vis at akselerasjonsvektoren $\vec a(t)$ er gitt ved
\[
\vec a(t) =
\begin{pmatrix}
-4\cos(2t) \\
-4\sin(2t) \\
0
\end{pmatrix}.
\]

\end{enumerate}
Det er nok ikke dumt å lagre denne GeoGebra-fila.

\section*{Kollisjon i 3D}

To partikkler beveger seg i rommet. Banene reperesenteres med følgende parameterframstillinger:

\[
a:
\begin{cases}
x(t)=\cos(2t) \\
y(t)=\sin(2t)-2 \\
z(t)=t-3
\end{cases},\, 0 \leq t \leq 5; \,\,\,\,\,\,\,\,\,\,\,\,\,\,\,\,\, b:
\begin{cases}
x(t)=t-2 \\
y(t)=\cos(2t) \\
z(t)=\sin(3t)
\end{cases},\, 0 \leq t \leq 5.
\]

\begin{enumerate}[label=\alph*)]
\item Legg både $a$ og $b$ inn i GeoGebra. 
\begin{lstlisting}
a = Kurve(cos(2t), sin(2t), t-3, t, 0, 5)
b = Kurve(t-2, cos(2t), sin(3t), t, 0, 5)
\end{lstlisting}

\item Legg til en glider for $t$. Legg deretter til punktene $A$ og $B$.
\begin{lstlisting}
A = a(t)
B = b(t)
\end{lstlisting}
Kolliderer partikklene?
\item Finn $\overrightarrow{AB}$. I GeoGebra kan du skrive:
\begin{lstlisting}
AB = Vektor(A,B)
\end{lstlisting}
Hva representerer denne vektoren?

\item Vis ved regning at $\overrightarrow{AB}$ er gitt ved
\[
\overrightarrow{AB} =
\begin{pmatrix}
t-2-\cos(2t) \\
\cos(2t) - \sin(2t)+2 \\
\sin(3t)-t+3
\end{pmatrix}.
\]

\item Finn $|\overrightarrow{AB}|$ og plott den i GeoGebra. Dette kan fint gjøres i CAS:
\begin{lstlisting}
L(x) := Vektor(b(x)-a(x))
f(x) := abs(L(x))
\end{lstlisting}
$x$ representerer nå tid, du kan ikke bruke $t$ siden den ble definert som en glider tidligere.

Basert på grafen til $f(x)$, kolliderer partikklene?

\item Gå utifra at partikklene er kuler og at begge har radius $r=0.5$. Når kolliderer de?\footnote{$t = 1.735...$.} I hvilke punkter er partikklene da?\footnote{$A(-0.946,-0.324,-1.26)$ og $B(-0.264,-0.946,-0.880)$.}
\end{enumerate}



\end{document}